\begin{table*}[!tb]
\newcommand{\factor}{}
\small
\centering
\caption{将 \method{} 与 \citet{tassa2018dmcontrol} 报告的无模型算法 A3C 和 D4PG 进行比较。对应训练曲线在其论文图~4 中以橙色线显示，在图~6 中以实线绿色线显示。我们据此估计 D4PG 达到 \method{} 最终性能所需的 episode 数，以估计数据效率增益。我们还加入使用真实模拟器而非学得动力学的 CEM 规划（$H=12,I=10,J=1000,K=100$），作为性能上界的估计。数字表示 5 个随机种子、10 条轨迹上的平均最终性能。}
\label{tab:control}
\begin{tabular}{llrrrrrrr}
\toprule
\textbf{方法} & \textbf{模式} & \textbf{诗集} & \textbf{\rot{Cartpole\\Swing Up}} & \textbf{\rot{Reacher\\Easy}} & \textbf{\rot{Cheetah\\Run}} & \textbf{\rot{Finger\\Spin}} & \textbf{\rot{Cup\\Catch}} & \textbf{\rot{Walker\\Walk}} \\ \midrule
A3C &自主性& 100,000 & 558 & 285 & 214 & 129 & 105 & 311 \\
D4PG &像素& 100,000 & 862 & 967 & 524 & 985 & 980 & 968 \\
\method (我们)&像素& 1,000 & 821 & 832 & 662 & 700 & 930 & 951 \\
CEM + 真模拟器&模拟状态& 0 & 850 & 964 & 656 & 825 & 993 & 994 \\
\midrule
\multicolumn{3}{l}{Data efficiency gain PlaNet over D4PG (factor)} & 250\factor & 40\factor & 500+\factor & 300\factor & 100\factor & 90\factor \\
\bottomrule
\end{tabular}
\end{table*}
